% Problem-section file following ../_template.tex.
\section{The quantum PCP conjecture}
\label{sec:problem-62}

\paragraph{Problem.}
Is the constant-relative-gap local Hamiltonian problem QMA-hard?  More
precisely, do there exist fixed integers $q,k\geq2$ and constants
$0\leq a<b\leq1$ such that the following promise problem is QMA-hard?  An
instance consists of $n$ subsystems of local dimension $q$ and
$m=\operatorname{poly}(n)$ positive semidefinite terms, each specified with
polynomially many bits, for which
\begin{equation}
  H:=\sum_{i=1}^{m}H_i,
  \qquad
  0\preceq H_i\preceq I,
  \qquad
  \lvert\operatorname{supp}(H_i)\rvert\leq k.
  \label{eq:p62-local-hamiltonian}
\end{equation}
Given the Hamiltonian in Eq.~\eqref{eq:p62-local-hamiltonian}, distinguish
the promised alternatives
\begin{equation}
  \lambda_{\min}(H)\leq am
  \qquad\text{and}\qquad
  \lambda_{\min}(H)\geq bm.
  \label{eq:p62-constant-gap-promise}
\end{equation}
Thus the gap in Eq.~\eqref{eq:p62-constant-gap-promise} is a fixed positive
fraction $(b-a)m$ of the number of local terms, independent of $n$.

\subsection*{Status}

Unsolved

\subsection*{Source}

Aharonov, Arad, and Vidick give the standard constant-relative-gap local
Hamiltonian formulation as Conjecture~1.3 of their quantum-PCP survey
\sourcecite{ref:p62-aharonov}{AAV13}.

\subsection*{Progress}

\begin{itemize}
  \item Aharonov, Arad, and Vidick review the QMA-completeness of local
  Hamiltonian with an inverse-polynomial promise gap and formulate the
  constant-gap strengthening in
  Eq.~\eqref{eq:p62-constant-gap-promise}
  \sourcecite{ref:p62-aharonov}{AAV13}.  Gap amplification methods known for
  classical constraint systems do not establish the quantum statement.

  \item Anshu, Breuckmann, and Nirkhe construct local Hamiltonians with the
  no-low-energy-trivial-states property from good quantum LDPC codes
  \sourcecite{ref:p62-anshu}{ABN23}.  This establishes the required
  low-energy entanglement phenomenon but does not prove QMA-hardness of the
  promise problem in Eqs.~\eqref{eq:p62-local-hamiltonian} and
  \eqref{eq:p62-constant-gap-promise}.

  \item Buhrman, Helsen, and Weggemans prove reductions among several
  quantum-PCP formulations and oracle separations showing that a proof of
  constant-gap local-Hamiltonian hardness must use nonrelativizing
  techniques \sourcecite{ref:p62-buhrman}{BHW25}.  Their results restrict
  possible proofs but neither prove nor refute the conjecture.
\end{itemize}

\subsection*{References}

\begin{enumerate}[label={},leftmargin=4.8em,labelsep=0.5em]
  \footnotesize
  \item[\textup{[AAV13]}]\label{ref:p62-aharonov}
  D. Aharonov, I. Arad, and T. Vidick,
  ``Guest Column: The Quantum PCP Conjecture,''
  \emph{ACM SIGACT News} \textbf{44}(2), 47--79 (2013).
  \href{https://doi.org/10.1145/2491533.2491549}%
       {doi:10.1145/2491533.2491549};
  \href{https://arxiv.org/abs/1309.7495}{arXiv:1309.7495}.

  \item[\textup{[ABN23]}]\label{ref:p62-anshu}
  A. Anshu, N. P. Breuckmann, and C. Nirkhe,
  ``NLTS Hamiltonians from Good Quantum Codes,'' in
  \emph{Proceedings of the 55th Annual ACM Symposium on Theory of Computing},
  1090--1096 (2023).
  \href{https://doi.org/10.1145/3564246.3585114}%
       {doi:10.1145/3564246.3585114};
  \href{https://arxiv.org/abs/2206.13228}{arXiv:2206.13228}.

  \item[\textup{[BHW25]}]\label{ref:p62-buhrman}
  H. Buhrman, J. Helsen, and J. Weggemans,
  ``Quantum PCPs: On Adaptivity, Multiple Provers and Reductions to Local
  Hamiltonians,'' \emph{Quantum} \textbf{9}, 1791 (2025).
  \href{https://doi.org/10.22331/q-2025-07-11-1791}%
       {doi:10.22331/q-2025-07-11-1791};
  \href{https://arxiv.org/abs/2403.04841}{arXiv:2403.04841}.
\end{enumerate}

\subsection*{Comment}

NLTS supplies a central structural prerequisite, and the known reductions
clarify which formulations and proof methods are viable, but neither gives
the constant-gap QMA-hardness required by
Eq.~\eqref{eq:p62-constant-gap-promise}.  No algorithmic obstruction that
would refute that hardness statement is known either.

\subsection*{Tag}

Quantum PCP conjecture; Hamiltonian complexity; Quantum complexity theory;
Quantum LDPC codes; Quantum locally testable codes

\subsection*{ID}

\texttt{op\_bca77ec42ddd1d5c}
